\begin{table}[ht]
\centering
\caption{Under-the-gun gap under alternative quantity controls.}
\label{tab:quantity_definition_robustness}
\begin{threeparttable}
\small
\setlength{\tabcolsep}{5pt}
\begin{tabular}{p{.42\linewidth}rrrr}
\toprule
Quantity control & Admin coef. & SE & Gap (\%) & $N$ \\
\midrule
No quantity control & -0.259 & 0.092 & 29.5 & 61,620 \\
Log accepted quantity & 0.145 & 0.133 & -13.5 & 61,620 \\
Raw accepted quantity / 1,000 & -0.243 & 0.086 & 27.5 & 61,620 \\
Within-item quantity percentile & -0.199 & 0.066 & 22.0 & 61,620 \\
\bottomrule
\end{tabular}
\begin{tablenotes}[flushleft]\footnotesize
\item \textit{Notes:} Each row estimates the administrative-versus-litigated urgent price specification with item, year, and PBU fixed effects. The coefficient is administrative minus litigated log negotiated price; the percentage column reports the implied litigated-over-administrative price gap. Quantity controls are included as listed. The BEC data used here do not contain a common dose-standardized quantity field, so the table varies the observed accepted-quantity transformation instead. Standard errors are clustered by PBU.
\end{tablenotes}
\end{threeparttable}
\end{table}
